\chapter{Conclusion}

Le travail presente dans ce rapport avait pour objectif d'etudier, de
reproduire et d'analyser de maniere critique l'approche \emph{Dial-In LLM}.
Cette demarche a conduit a l'etablissement d'une baseline de reproduction
documentee, a l'identification de plusieurs limites de reproductibilite et a
l'evaluation de modifications methodologiques destinees a ameliorer le
pipeline.

\section{Resultats principaux}

Le premier resultat du projet est la construction d'une baseline de
reproduction defendable sur les quatre benchmarks anglais du papier. Cette
reproduction n'est pas strictement equivalente au protocole original, en raison
de l'absence des checkpoints fine-tunes utilises par les auteurs, mais elle est
suffisamment alignee pour rendre possible une comparaison quantitative
pertinente.

Le deuxieme resultat concerne les evaluations complementaires. Plusieurs
propositions ont ete examinees, mais la plus convaincante est la politique
adaptative \texttt{sqrt} pour la selection de \texttt{K}. Cette modification
ameliore sensiblement la qualite du clustering sur les benchmarks les plus
sensibles et reduit simultanement la fragmentation finale.

Le troisieme resultat est d'ordre analytique. Le projet a permis de distinguer
clairement trois niveaux de comparaison :
\begin{itemize}[leftmargin=*]
  \item les resultats rapportes dans le papier ;
  \item la baseline de reproduction etablie dans notre environnement ;
  \item la meilleure configuration amelioree issue des ablations.
\end{itemize}

Cette structuration est essentielle, car elle permet de separer rigoureusement
ce qui releve de la reproduction et ce qui constitue une contribution
experimentale propre.

\section{Interpretation generale}

L'analyse d'ensemble conduit a une conclusion nuancee. La reproduction seule se
revele deja competitive sur plusieurs benchmarks. La configuration amelioree
surpasse ensuite la reproduction sur trois benchmarks sur quatre et reduit une
partie du decalage avec le papier, tout en depassant legerement le score de
reference sur MTOP(I). En revanche, aucune configuration unique ne domine
systematiquement sur l'ensemble des cas etudies, comme le montre le comportement
de Bank77.

Cette heterogeneite constitue un resultat en soi. Elle indique que les
performances de la methode restent fortement dependantes de la structure des
benchmarks et du comportement de certaines composantes du pipeline, en
particulier la selection de \texttt{K}, la stabilite du jugement de coherence
et le niveau de fragmentation final.

\section{Limites du travail}

Plusieurs limites doivent etre explicitees.
\begin{itemize}[leftmargin=*]
  \item Les checkpoints fine-tunes utilises par les auteurs pour les benchmarks
  anglais ne sont pas publiquement disponibles.
  \item Le recours a \texttt{Mistral-7B-Instruct} constitue une approximation
  operationnelle, mais non equivalente, du \texttt{Llama-7B} fine-tune du
  papier.
  \item Toutes les evaluations complementaires n'ont pas ete conduites de
  maniere exhaustive sur les quatre benchmarks.
  \item Les performances demeurent sensibles a certains choix de configuration,
  ce qui limite la portee d'une lecture uniforme des resultats.
\end{itemize}

Ces limites ne remettent pas en cause l'interet du travail, mais elles bornent
explicitement son domaine de validite et permettent de distinguer une
reproduction approchee d'une replication stricte.

\section{Apport scientifique du projet}

L'apport principal de ce travail tient a l'articulation entre reproduction,
analyse critique et proposition methodologique. Il montre :
\begin{itemize}[leftmargin=*]
  \item qu'il est possible de reconstruire le coeur de \emph{Dial-In LLM} dans
  un environnement public et controle ;
  \item que les ecarts au papier peuvent etre interpretes a partir de
  composantes techniques identifiables ;
  \item que certaines modifications, notamment sur la gestion de \texttt{K},
  apportent un gain mesurable ;
  \item que l'evaluation de la methode doit rester benchmark-dependante et ne
  peut etre reduite a un commentaire global sur les scores.
\end{itemize}

En ce sens, le projet satisfait les attendus d'un travail tutore en
apprentissage automatique et en science des donnees : appropriation d'un etat
de l'art, reproduction critique d'une methode recente, evaluations
complementaires et formulation de propositions argumentees.

\section{Ouverture}

La phase experimentale peut desormais etre consideree comme stabilisee. La
suite du travail releve principalement de la consolidation editoriale et de la
valorisation des resultats : finalisation du rapport PDF, preparation des
supports de restitution et presentation synthetique des contributions du groupe
dans un cadre academique.
